\subsection{Buying capacity when models stay resident}
\label{exec17:s8-residency}

The production-instance description has a further implication that a throughput
table misses. Requests can revisit an already loaded model, while retaining
that model consumes a slot between requests. An acquisition decision therefore
depends on the model-retention policy as well as the number of licenses. The
Bayesia overview explicitly allows sequential loading and unloading to reduce
the simultaneous-instance requirement. It does not specify the latency of
these operations, a caching policy, or a guarantee that one model instance can
serve concurrent requests. Those are quantities to measure. The following
reference isolates residency while keeping actual product behavior unresolved
\citep{exec13_bayesia_api,exec17_s8_residency}.

Consider a single worker receiving requests for immutable model versions from
a finite set. Loading a model consumes five abstract time units, inference
consumes one, and eviction is instantaneous. No model is loaded initially.
Requests are processed in their declared order, with no prefetch and no overlap
of loading and inference. The bank buys capacity for $m$ resident models,
including idle ones. These assumptions deliberately exclude parallel speedup:
the question is whether buying an additional slot saves reloads. A model name
includes its parameter and schema version; a revised balance-response network
is a different object even when an analyst gives it the same display name.

Write the ordered model requests as $x_1,\ldots,x_K$, and let $S$ be the
resident set just before a request. A hit occurs if $x_t\in S$. A miss requires
loading $x_t$ and, if $|S|=m$, selecting an eviction. For a known batch, let
$J_t(S)$ denote the smallest number of remaining loads from time $t$. Retaining
unused models is free apart from capacity, so an optimum can delay an eviction
until the slot is needed. The recurrence is
\begin{equation}
\begin{split}
 J_{K+1}(S)&=0,\\
 J_t(S)&=J_{t+1}(S),\qquad x_t\in S,\\
 J_t(S)&=1+\min_{S'\in\mathcal E_m(S,x_t)}J_{t+1}(S'),\qquad x_t\notin S,
\end{split}
\label{exec17:s8-resident-bellman}
\end{equation}
where $\mathcal E_m(S,x)$ is $\{S\cup\{x\}\}$ if $|S|<m$, and is
$\{(S\setminus\{y\})\cup\{x\}:y\in S\}$ if $|S|=m$.
Thus every candidate next set contains the requested model and respects the
purchased capacity. The recurrence is exact by conditioning on the next
eviction and then optimizing the remaining suffix. It is not an estimator of
a future workload. Under the stated serial model, a schedule with $N$ loads
has completion time
\begin{equation}
 T=K+5N.
 \label{exec17:s8-resident-time}
\end{equation}
Changing the load or inference duration replaces those coefficients; overlapping
operations or concurrent calls require a different state and time model.

Take the six-request batch $(A,B,C,A,B,C)$. One slot requires six loads and
finishes at 36. Three slots load each distinct model once and finish at 21.
With two slots, the following resident sets attain four loads:
\begin{equation}
 \{A\},\quad\{A,B\},\quad\{A,C\},\quad\{A,C\},\quad\{B,C\},\quad\{B,C\}.
 \label{exec17:s8-resident-witness}
\end{equation}
Load at requests 1, 2, 3 and 5; requests 4 and 6 are hits. Completion time is
26. At least four loads are necessary: the first three distinct requests each
require a load, and loading $C$ in the third request with capacity two evicts
either $A$ or $B$. Both are requested again, so at least one further load is
unavoidable. This lower bound and the witness establish the exact two-slot
optimum. A separate enumeration of all feasible resident-set paths gives the
same result, including paths that voluntarily evict early
\citep{exec17_s8_residency}.

The offline optimum needs information about the remaining batch. A familiar
online rule evicts the least recently used resident model. For the same
two-slot sequence its resident sets are $\{A\},\{A,B\},\{B,C\},\{C,A\},
\{A,B\},\{B,C\}$. Every request is a miss, so it takes 36, exactly the
one-slot time. The extra license has supplied storage capacity but this policy
has not converted it into a time saving. The distinction is commercially
material: a demonstration that knows the complete batch cannot support a
streaming service-level promise for an unspecified eviction policy.

Reordering requests to group identical models would answer a different
scheduling problem. It may be legitimate for independent overnight reports;
it may be impossible when the next servicing request depends on an earlier
decision or when each request has a deadline. The acquisition case must say
whether such reordering is permitted. Likewise, preloading all three models
before starting the stopwatch both consumes capacity and incurs work. Excluding
that preparation would favor the larger resident configuration for the wrong
reason. The reference starts with an empty resident set so the full preparation
cost is visible.

\subsection{A workload law and a deadline produce different purchases}
\label{exec17:s8-resident-workload}

An expected-cost quote needs a workload law. Suppose, only for this example,
that each of the six requests independently chooses $A$, $B$ or $C$ with
equal probability. There are $3^6=729$ equally likely batches. This is an
authored finite law, not an inferred distribution of bank demand. It permits
both complete enumeration and an independent probability calculation for the
online least-recently-used rule.

For one slot, the first request always loads and each later request differs
from its predecessor with probability $2/3$. For three slots, loads equal the
number of distinct requested models. Therefore
\begin{equation}
 \E N_1=1+5(2/3)=\frac{13}{3},\qquad
 \E N_3=3\{1-(2/3)^6\}=\frac{665}{243}.
 \label{exec17:s8-resident-means13}
\end{equation}
With two slots, before request $t\geq2$ there is only one resident model
exactly when all the preceding requests are equal. That event has probability
$3^{2-t}$. Conditional on it, the next request misses with probability $2/3$;
otherwise there are two resident models and the next independent request
misses with probability $1/3$. Hence
\begin{equation}
 \E N_2=1+\sum_{t=2}^{6}\left\{\frac13+\frac13\,3^{2-t}\right\}
       =\frac{769}{243}.
 \label{exec17:s8-resident-mean2}
\end{equation}
The independence of the next request from the resident state is essential to
this calculation. Changing model frequencies or making the next model depend
on previous outputs changes the law. Exhaustive evaluation of all 729 batches
matches these three rational expectations exactly; it is not Monte Carlo
evidence with an estimated standard error \citep{exec17_s8_residency}.

Now assign an illustrative allocated capacity charge of USD~0, 120 and 300
per six-request planning batch to one, two and three slots. These are
incremental accounting assumptions, not a vendor's billing units or quotes.
Price a time unit at USD~20, including the stipulated common preparation and
waiting cost. For a known batch, the objective is the charge plus $20T_m$;
for the specified random online workload, it is the charge plus $20\E T_m$.
The comparison is
\begin{equation}
\begin{array}{c|rrr}
 \text{resident slots}&1&2&3\\ \hline
 \text{known-batch optimum time}&36&26&21\\
 \text{known-batch total USD}&720&640&720\\
 \text{online mean total USD}&1660/3&135220/243&168560/243
\end{array}.
\label{exec17:s8-resident-prices}
\end{equation}
Thus the known batch favors two slots; the random online workload favors one,
with expected cost approximately USD~553.33 against USD~556.46 for two and
USD~693.66 for three. The small expected difference between one and two is a
property of these supplied numbers. It would be irresponsible to regard it
as a measured procurement margin without observing the workload and costs.

A hard deadline is a separate constraint. Over the finite support of six-request
sequences, one slot and two-slot least-recently-used scheduling both have
worst completion time 36. Three slots have worst completion time 21, because
there are at most three distinct models. A requirement to finish every supported
sequence within 30 units therefore selects three among these online
configurations, despite its higher mean cost. For the known batch, the exact
two-slot plan finishes at 26 and is eligible. The facts available before
purchase decide whether that cheaper plan is a relevant comparator.

This is also a limitation of any deterministic online two-slot policy under
an unrestricted adaptive arrival model. An adversary observing the resident
set can always request a missing one of the three models, forcing a load at
each of the six requests. That argument does not assert that bank arrivals are
adversarial. It shows why a worst-case claim cannot be inferred from an average
or a prescient batch schedule. If requests are random and independent, use
that declared law. If the bank requires a deadline only with a specified
probability, estimate that tail under an appropriate workload model and
report the failed fraction. Do not silently reinterpret ``every request'' as
``the average request.''

The example has deliberately held scientific acceptance fixed across slots:
all model versions and outputs are identical and no inference is skipped.
An apparent speedup obtained by retaining an old model after an approved
parameter update changes that condition. If $A$ becomes $A'$ midway through
a batch, the residency identity and load count must change. If the policy
changes the model selected after observing an output, future requests are
endogenous to the earlier computation and the independent uniform calculation
no longer applies. A prototype should retain those version and decision events
so that the capacity estimate corresponds to the actual approved analysis.

\subsection{What the capacity prototype must reveal before a quote}

The loaded-instance model identifies a small set of observations with direct
acquisition consequences. For each permitted configuration, measure loading,
license validation, readiness, inference and unload times separately; keep
resident versions and model-file locations throughout the trace. An engine
that runs quickly after loading can still be unsuitable for frequent model
changes. A route with customer-resident model files can still depend on the
documented license-server connection. An interrupted or unlicensed load must
remain in the assigned workload, with its actual fallback or unresolved
deadline; it is not a zero-time successful request
\citep{exec13_bayesia_api}.

Two comparisons should then be reported. A known-batch reference measures
what a sufficiently informed scheduler could attain under the same capacity
and ordering constraints. A deployable online policy measures what the
proposed service can attain using only information available at each request.
The gap between them is a scheduling opportunity, not evidence that the vendor
engine is inefficient. Original loaded-model traces allow the bank to ask
whether an additional instance or a better retention policy would address the
observed bottleneck. A mean CPU time without such traces cannot distinguish
those purchases.

License and runtime failures require a further state. The finite time formula
assumes every required load is permitted and finishes in its stipulated time.
If validation is unavailable with no known upper delay, this model supplies
no finite deadline guarantee. The response is to keep that operational input
open, evaluate an allowed fallback and avoid quoting unsupported availability.
The same separation applies to a SMILE search interrupted during an iteration:
extra licensed capacity does not turn an incomplete graph into a completed
scientific result. A capacity proposal should preserve the completed/incomplete
state while measuring its operational cost \citep{exec13_smile_pc}.

Finally, allocate recurring capacity charges to a planning workload only once.
A license shared by funding and servicing analyses is a common fixed resource;
charging its full annual fee independently to each journey overstates total
cost, while omitting one journey's occupancy understates contention. The
prototype should report joint resident intervals and the rule for allocating
their expense. A price negotiation can then use both a scientifically matched
output requirement and a concrete instance quantity. The proposal's immediate
deliverable is this testable purchasing model. Actual durations, permitted
concurrency, deployment rights, prices and bank workload remain observations
to obtain after the narrow prototype shows merit.
